% 一维阈值与单区间的VC维：可打散的点集与多一个点时的反例。
\begin{tikzpicture}[figure picture,foundation picture,x=1mm,y=1mm,font=\figurefont\footnotesize]
  \fill[FigurePaper] (0,0) rectangle (108,74);
  \node[font=\figurefont\footnotesize,text=FigureMuted] at (11,66) {假设空间};
  \node[font=\figurefont\footnotesize,text=FigureMuted] at (47,66) {能打散的点集};
  \node[font=\figurefont\footnotesize,text=FigureMuted] at (90,66) {多一个点时的反例};
  \draw[FigureGrid] (4,60) -- (104,60);
  \draw[FigureGrid] (4,34) -- (104,34);
  \draw[FigureGrid] (22,5) -- (22,70);
  \draw[FigureGrid] (73,5) -- (73,70);

  \node[font=\figurefont\footnotesize,text=FigureInk] at (11,50) {阈值};
  \node[text=FigureInk] at (11,41) {$\operatorname{VC}=1$};
  \draw[FigureRed,line width=.75pt,fill=FigurePaper] (37,49) circle (1.8);
  \fill[FigureBlue] (57,49) circle (1.8);
  \node at (37,40) {$(0)$};
  \node at (57,40) {$(1)$};
  \draw[FigureMuted,line width=.5pt] (79,48) -- (101,48);
  \fill[FigureBlue] (83,48) circle (1.8);
  \draw[FigureRed,line width=.75pt,fill=FigurePaper] (97,48) circle (1.8);
  \draw[LancetGreen,line width=.85pt,->] (83,54) -- (101,54);
  \node at (90,40) {$(1,0)$无法实现};

  \node[font=\figurefont\footnotesize,text=FigureInk] at (11,24) {区间};
  \node[text=FigureInk] at (11,15) {$\operatorname{VC}=2$};
  \foreach \cx/\a/\b in {29/0/0,41/0/1,53/1/0,65/1/1}{
    \draw[FigureMuted,line width=.5pt] (\cx-4,21) -- (\cx+4,21);
    \ifnum\a=1
      \fill[FigureBlue] (\cx-2,21) circle (1.25);
    \else
      \draw[FigureRed,line width=.75pt,fill=FigurePaper] (\cx-2,21) circle (1.25);
    \fi
    \ifnum\b=1
      \fill[FigureBlue] (\cx+2,21) circle (1.25);
    \else
      \draw[FigureRed,line width=.75pt,fill=FigurePaper] (\cx+2,21) circle (1.25);
    \fi
    \node at (\cx,12) {$(\a,\b)$};
  }
  \draw[FigureMuted,line width=.5pt] (77,21) -- (103,21);
  \fill[FigureBlue] (80,21) circle (1.6);
  \draw[FigureRed,line width=.75pt,fill=FigurePaper] (90,21) circle (1.6);
  \fill[FigureBlue] (100,21) circle (1.6);
  \draw[LancetGreen,line width=.85pt] (80,28) -- (100,28);
  \draw[LancetGreen,line width=.85pt] (80,26.5) -- (80,29.5);
  \draw[LancetGreen,line width=.85pt] (100,26.5) -- (100,29.5);
  \node at (90,12) {$(1,0,1)$无法实现};
\end{tikzpicture}
